% Appendix A -- The shared compiler corpus.  Source (Law 1): os/tps/
% CORPUS.md s2 (the controlled register: stem, DOS name, family, coverage)
% and the Makefile variable TPS_COMPILER_CORPUS (the order, shared by the
% host and InitechDOS legs).  The table is checked against the shipped
% corpus through the replayed run of TIN-P-03 (make test-compiler-os prints
% one PASS line per program, in this order); a program added, removed or
% renamed changes that run and turns test-pubs-turbo red.
\Gates{test-compiler, test-compiler-os, test-pubs-turbo}
\chapter{The shared compiler corpus}
\ix{corpus}

Every build of Turbo Initech must compile these seventeen programs, and the
code it makes for each must print exactly what the code made by the Free
Pascal compiler prints, which in turn must match the output worked out by
hand for that program.  The programs are in \file{os/tps/fixtures/}; each
has a sibling file \file{.golden} holding the expected output.  On
InitechDOS they run under the names in the second column.

\begin{BorTable}{The corpus}{p{39mm}p{17mm}p{40mm}}
\Hd{Program} & \Hd{Name} & \Hd{What it exercises}\EndHd
\code{gen\_tiny} & \file{CTINY} & loops, a branch, two value parameters, a
  local\\
\code{gen\_bool\_deep} & \file{CBOOL} & every relation, \rw{and}, \rw{or},
  \rw{not}, precedence\\
\code{gen\_control\_deep} & \file{CCTRL} & \rw{if}, \rw{while}, \rw{for} in
  both directions, \rw{repeat}\\
\code{gen\_char\_deep} & \file{CCHAR} & constants, characters, \std{ord}{Ord},
  \std{chr}{Chr}\\
\code{gen\_func\_shared} & \file{CFUNCS} & recursion, mutual recursion with
  \rw{forward}, a \rw{var} swap\\
\code{gen\_func\_deep} & \file{CFUNC} & locals, value and \rw{var}
  parameters, aliasing\\
\code{gen\_recursion\_depth} & \file{CDEPTH} & negative \rw{div} and \rw{mod},
  thirteen recursive frames, complete evaluation\\
\code{gen\_array\_shared} & \file{CARRAYS} & global and local arrays, bounds,
  elements as \rw{var} arguments\\
\code{gen\_array\_deep} & \file{CARRAY} & integer, Boolean and character
  arrays; a local array\\
\code{gen\_record\_shared} & \file{CRECS} & fields, record \rw{var}
  parameters and copies, arrays of records\\
\code{gen\_record\_deep} & \file{CRECORD} & record layout, local records,
  nested references\\
\code{gen\_local\_record\_array} & \file{CLRECORD} & a local array of
  records with a negative lower bound\\
\code{gen\_string\_shared} & \file{CSTRS} & concatenation, length, indexing,
  comparison\\
\code{gen\_string\_deep} & \file{CSTRING} & truncation, every comparison,
  recursion with strings\\
\code{gen\_string\_temporaries} & \file{CSTRTEMP} & four nested
  temporary strings; doubled apostrophes\\
\code{gen\_file\_shared} & \file{CFILES} & writing and reading back bytes and
  blocks; a local file\\
\code{gen\_fileio\_deep} & \file{CFILEIO} & the same file operations, written
  independently\\
\end{BorTable}

The run of procedure TIN-P-03 in Chapter 4 reports each program in this
order, first on the development system and then on InitechDOS.
